\section{强化学习算法设计}

\subsection{PRIME：从Value Model视角出发}

崔干曲团队在2024年初（DeepSeek R1之前）试图复刻O1时，面临的核心问题是：\textbf{如何突破Sparse Reward导致的Credit Assignment困难？}

传统方法（如GRPO）仅使用最终的Verifiable Reward作为训练信号，这意味着模型只能从最终结果获得反馈，无法精确知道推理过程中的哪些步骤是正确的。PRIME的核心思路是引入\textbf{Implicit PRM}（隐式过程奖励模型），通过Log Probability的方式自动进行Credit Assignment。

\begin{importantbox}{PRIME的核心设计思想}
\begin{itemize}
    \item \textbf{问题}：PPO的Value Model在大语言模型上训练初期预测极差，导致reward先下降再回升（``dip-then-recover''现象），引入额外的训练不稳定性
    \item \textbf{方案}：用Log Probability建模Value，而非传统的Linear Head方式。训练前后，Reward Model仍然是语言模型形式，避免了丢弃embedding层接linear head的粗暴做法
    \item \textbf{优势}：这是一种更自然、更健壮的Value建模方式，尤其在训练初期不会出现严重的Value prediction偏差
\end{itemize}
\end{importantbox}

从今天的视角回看，PRIME其实是一种\textbf{Value Model的替代建模}方式。虽然当前GRPO等无Value Model的方法已经能训得很好，但崔干曲认为，在未来更复杂的场景（如Agentic RL中存在内部feedback的情况），如何将多源feedback转化为Process Reward或Value Estimation，仍然是一个值得深入研究的问题。

此外，PRIME与Uncertainty Distillation方法有深层的联系——两者都在最小化Policy Model与Teacher/Reward Model之间的KL散度，本质上是一种将Reward Model知识蒸馏到Policy Model的过程。

\begin{knowledgebox}{Value-based Actor-Critic方法的未来}
当前GRPO等方法因简单高效而大行其道，但Value-based方法并非过时。崔干曲预测，随着任务复杂度提升（更多的feedback来源、更长的交互链路），社区将重新关注Value-based Actor-Critic方法。这一方向仍有大量可深挖的空间。
\end{knowledgebox}

\subsection{Reinforce++：回归经典的智慧}

胡建设计Reinforce++的出发点带有"回归经典"的色彩。在传统Multi-Agent RL时代，PPO在大规模游戏场景中表现非常稳定。来到大模型时代后，胡建的直觉是：经典方法中的许多细节设计（如Global Batch Norm）是经过前人深思熟虑的，应该被继承而非丢弃。

\begin{importantbox}{Reinforce++的设计哲学}
\begin{enumerate}
    \item 从PPO出发，\textbf{去掉Critic}（避免大模型上Value Model的训练困难）
    \item \textbf{保留PPO中经过验证的工程技巧}：Global Batch Norm等
    \item 验证了经典技巧在新场景下仍然有效，"长期存在的东西一定有它存在的道理"
\end{enumerate}
\end{importantbox}

胡建进一步提出了一个重要观点：未来的RL算法不再是一个单独的算法，而是\textbf{一组优良Recipe的组合}。例如：Truncated Importance Sampling + On-Policy Learning + Batch Norm + 合适的Infra实现，共同构成一套稳定的训练方案。

\subsection{GSPO：Sequence Level优化的回归}

郑楚杰团队在训练千问2507版本时发现，GRPO在后期总会出现稳定性问题。他们怀疑根本原因在于\textbf{优化目标本身}：GRPO的优化目标本应是Sequence Level的，但与Token Level的优化目标之间缺乏直接的理论关联。

\begin{importantbox}{GSPO的核心创新}
\begin{itemize}
    \item \textbf{回归第一性原理}：直接推导Sequence Level的优化目标
    \item \textbf{去掉KL Penalty}：在235B及更大模型上，去掉KL Penalty后训练速度显著提升，且GSPO能保持稳定训练
    \item \textbf{实际应用}：GSPO被应用于千问2507发版以及后续模型（如Consonant X）的训练
\end{itemize}
\end{importantbox}

\begin{warningbox}{GSPO的局限性}
初版GSPO没有考虑训推不一致（Train-Inference Mismatch）问题。后续在内部迭代中增加了对Mismatch的修复，产生了GSPO-VR等改进算法，进一步提升了稳定性。
\end{warningbox}

\subsection{从第一性原理出发的算法设计}

李英儒从强化学习理论研究者的视角出发，强调了两个核心观点：

\textbf{第一，回归第一性原理。}许多当前的训练崩溃问题源于对基本概念的忽视。例如，很多框架在Recompute阶段将重新计算的概率当作Rollout概率使用，这不符合直觉——真正产生采样的概率应来自推理引擎（如vLLM），而非训练引擎的Recompute。

\textbf{第二，理论应为实践服务。}李英儒认为，做理论不是为了复杂而复杂，而是为了：(1) 提供简单的解释；(2) 为设计更简单有效的Scalable算法提供方法论指导。

\begin{knowledgebox}{Off-Policy Issue的历史渊源}
Train-Inference Mismatch问题本质上是经典的Off-Policy Issue在大模型时代的新表现。这一问题在传统RL中已经研究了数十年。李英儒自2019年起就在游戏RL中研究类似问题（Divergence of Entity），当时做的工作与现在的KL Penalty方法非常相似。从TRPO（2015年）到现在的各种IS修正方法，核心都是关于\textbf{单调性能提升的保证}。
\end{knowledgebox}

\subsection{"好的RL算法"应具备什么特点？}

四位嘉宾达成了高度一致的共识：

\begin{importantbox}{好的RL算法（Recipe）的标准}
\begin{enumerate}
    \item \textbf{简单}：每个Trick应易于实现，新手看一眼就能理解其作用
    \item \textbf{符合直觉}：设计应有清晰的motivation，不是黑箱式的"加了就有效"
    \item \textbf{有效}：能在各种模型规模和任务上稳定训练，并产生真实的性能收益
    \item \textbf{组合性}：好的"算法"实际上是一套Recipe——基本优化目标（如Reinforce + IS Correction）加上一系列经过验证的Trick（PPO Clip、Batch Norm、Language Consistency Penalty等）
\end{enumerate}
\end{importantbox}

郑楚杰进一步指出：与其说"好的RL算法"，不如说"好的RL配方（Recipe）"。现实中，大家摸索出一套稳定的配方后就开始使用，并不会给这套Trick的组合专门起个名字。

\subsection{本章小结}

RL算法设计正在从追求单一突破性算法转向追求一套稳定、简洁、可组合的Recipe。PRIME代表了Value Model建模的新思路，Reinforce++体现了回归经典的智慧，GSPO展示了第一性原理推导的力量。四位嘉宾一致认为，未来的RL算法创新将更多来自Infra与算法的联合优化，而非纯粹的理论推导。


